\chapter{Analyse des besoins et conception du système}

\minitoc

\section{Besoins fonctionnels}
% TODO: lister les fonctions que le système doit assurer : liste à puces, tableau
% des acteurs et des cas d'utilisation ou diagrammes, selon les besoins du projet.
[Listez les besoins fonctionnels.]

\section{Besoins non fonctionnels}
% TODO: performances, sécurité, ergonomie, évolutivité, maintenabilité, etc.
[Listez les besoins non fonctionnels.]

\section{Exigences de sécurité}
% Section facultative -- à conserver si la sécurité ou la conformité faisait partie du projet
% (par exemple, renforcement de la sécurité selon l'OWASP ou respect d'une exigence ISO).
[Décrivez les exigences de sécurité ou de conformité, ou supprimez cette section.]

\section{Carnet de produit}
% TODO: présenter le carnet de produit et la répartition des tâches, généralement sous forme de tableau :
% Identifiant | Fonctionnalité majeure | Récit utilisateur | Priorité
[Insérez le carnet de produit, par exemple sous forme de tableau ou de liste de fonctionnalités majeures et de récits utilisateurs.]

\subsection{Priorisation du carnet de produit}
% TODO: préciser la méthode de priorisation (MoSCoW, WSJF, points d'effort, etc.) -- facultatif.
[Décrivez la méthode de priorisation, si nécessaire.]

\section{Architecture du système}

\subsection{Architecture physique}
% TODO: présenter le déploiement -- serveurs, conteneurs, réseau. Ajouter un schéma si disponible.
[Décrivez l'architecture physique et le déploiement.]

\subsection{Architecture logique}
% TODO: présenter les couches, les modules et leurs interactions.
[Décrivez l'architecture logique.]

\subsection{Présentation des composants}
% TODO: décrire brièvement chaque composant ou module principal.
[Décrivez les principaux composants.]

\section{Conclusion}
% TODO: conclure le chapitre en 2 à 3 phrases.
[Rédigez une brève conclusion du chapitre.]
